% ID: SEL-D0014-BA-E01
% ESTADO: cotejada. Segunda pasada visual completa y borrador compilado revisados.
% Fuente: Selectividad/matematicas_aplicadas/matematicas_aplicadas_P_B.pdf
% SHA-256: 51601c4f5a6c1dfcd860f621a659f8b9c09f53ff4cd31a1fa1e827dbd8f06159
% Página: 1; bloque A; ejercicio original 1.

\textbf{(2.5 puntos)} Con el fin de recaudar dinero para el viaje de fin de
curso, los alumnos de un instituto van a poner a la venta dos tipos de bolsas
de merienda. El primer tipo contendrá dos bocadillos, un refresco y una pieza
de fruta y el segundo tipo tendrá un bocadillo, un refresco y dos piezas de
fruta. Por cada bolsa del primer tipo cobrarán 6 euros y por las del segundo
tipo 5 euros. Sabiendo que disponen de 120 bocadillos, 70 refrescos y 110
piezas de fruta y que se tiene garantizada la venta de todas las bolsas,
\textquestiondown{}cuántas convendría preparar de cada tipo para que la
cantidad de dinero obtenida por su venta sea máxima y a cuánto asciende la misma?

\textquestiondown{}Es posible que vendan 40 bolsas de cada tipo?
\textquestiondown{}Hay alguna posibilidad de que el importe de las ventas sea
de 410 euros?
